% 配置字体、颜色、框（tcolorbox）样式

% 配置中文主字体：指定fonts目录下的Fandol字体
\usepackage{fontspec}

% 中文主字体：Fandol 宋体（含粗体、斜体配置）
\setCJKmainfont[
  Path       = fonts/,    % 字体文件所在相对路径（确保路径正确）
  Extension  = .otf,      % 字体文件扩展名
  BoldFont   = FandolHei-Regular,  % 粗体字体名
  ItalicFont = FandolKai-Regular   % 斜体（楷体）字体名
]{FandolSong-Regular}

% 中文无衬线字体：Fandol 黑体
\setCJKsansfont[
  Path       = fonts/,
  Extension  = .otf,
  BoldFont   = FandolHei-Bold
]{FandolHei-Regular}

% 中文等宽字体：Fandol 仿宋
\setCJKmonofont[
  Path       = fonts/,
  Extension  = .otf
]{FandolFang-Regular}

% 中文字族（同样修正参数格式）
\setCJKfamilyfont{zhsong}[
  Path       = fonts/,
  Extension  = .otf,
  BoldFont   = FandolHei-Regular
]{FandolSong-Regular}

\setCJKfamilyfont{zhhei}[
  Path       = fonts/,
  Extension  = .otf,
  BoldFont   = FandolHei-Regular
]{FandolHei-Regular}

\setCJKfamilyfont{zhfs}[
  Path       = fonts/,
  Extension  = .otf
]{FandolFang-Regular}

\setCJKfamilyfont{zhkai}[
  Path       = fonts/,
  Extension  = .otf
]{FandolKai-Regular}

% 字体命令
\renewcommand{\songti}{\CJKfamily{zhsong}}
\renewcommand{\heiti}{\CJKfamily{zhhei}}
\renewcommand{\fangsong}{\CJKfamily{zhfs}}
\renewcommand{\kaishu}{\CJKfamily{zhkai}}

% color settings
\definecolor{shallow-frame}{RGB}{211, 211, 211}
\definecolor{shallow-background}{RGB}{240, 240, 240}
\definecolor{warm-background}{RGB}{239, 239, 221}
\definecolor{blue-background}{RGB}{214, 220, 229}
\definecolor{deepblue-background}{RGB}{141, 170, 212}
\definecolor{darkgray-background}{RGB}{217, 217, 217}
\definecolor{mydarkgray}{RGB}{76, 76, 76}

\usepackage[most]{tcolorbox}

% tcolorbox
\newtcolorbox{mybox}[1]{
  colback=darkgray-background!40,
  coltitle=black,
  colframe=darkgray-background,
  colbacktitle=darkgray-background!60,
  fonttitle=\heiti,
  title={#1}
}

\newtcolorbox{mynotice}{
  colback=darkgray-background!40,
  colframe=darkgray-background,
}

\newtcolorbox{mynote}[1]{
  title={#1},
  coltitle=black,
  fonttitle=\heiti,
  colback=white,
  colframe=darkgray-background!60,
}